% ----------------------------------------------------------------------------------
% Chapter 3: Methodology
% ----------------------------------------------------------------------------------
% Describe the theoretical approach, algorithms used (e.g., Cartographer, Nav2, MoveIt2), and experimental design.
% Maintain an objective, third-person tone (or passive voice as appropriate).

\section{System Overview}
The proposed system relies on a tightly coupled software architecture leveraging ROS 2...

\section{Drive Mode Selection: Differential Drive over Ackermann Steering}
The LIMO PRO platform supports two drive modes: Ackermann steering (front-wheel steering
with a fixed wheelbase turn radius) and differential drive (independent left--right wheel
velocity commands). The system was initially prototyped in Ackermann mode, but a
fundamental constraint of Ackermann kinematics made it unsuitable: a minimum-radius turn
requires forward motion, which means the vehicle cannot rotate in place. For an
approaching task---navigate to a box, stop in front, then align the gripper precisely---the
inability to perform in-place rotation forces the planner into multi-arc arcing manoeuvres
that are sensitive to odometry drift and make the final alignment geometry complex. The
differential-drive mode removes this constraint entirely: the robot can spin in place to
any heading after stopping, and the Nav2 \texttt{RotateInPlace} behaviour is natively
supported. Additionally, the differential-drive odometry model requires only two
parameters (wheel separation and wheel radius) and is directly verifiable against the
physical platform's measured geometry, whereas Ackermann odometry involves a steering
angle measurement that adds a calibration dependency. For these reasons, differential
drive was adopted as the sole operational mode for all experiments in this work.

\section{Mapping and Localization (SLAM)}
Environment mapping is performed using SLAM Toolbox~\cite{macenski2021slam} in
synchronous mapping mode, fusing LiDAR range data with wheel odometry to construct a
2D occupancy grid of the workspace. The EAI T-mini Pro LiDAR is configured to a
forward arc of $\pm120^{\circ}$ (2.094\,rad half-angle) with a minimum valid range
of 0.10\,m. Both constraints are physically motivated: the sensor's angular subtended
arc at the chassis sides begins at approximately $\pm108^{\circ}$, so beams beyond
that angle strike the robot's own body and produce spurious occupied cells; the 0.10\,m
minimum eliminates the residual chassis returns that lie within the sensor's field of
view. Operating with the unrestricted $360^{\circ}$ field of view baked 55 isolated
noise pixels into the map as phantom obstacle cells, degrading localisation and
inflating the global costmap. Once the sensor parameters were corrected, the mapping
session was repeated from scratch.

Localisation during navigation is performed by the Adaptive Monte Carlo Localisation
(AMCL) filter, which matches the live scan against the pre-built map to maintain a
particle-based pose estimate. The AMCL laser minimum range is set to match the physical
sensor threshold (0.10\,m), ensuring that AMCL and the sensor model agree on which
returns to treat as valid.

\section{Autonomous Navigation}
Autonomous navigation is realised with the Nav2 framework~\cite{macenski2020nav2}
configured for a differential-drive base. The choice of differential drive over the
Ackermann steering mode available on the LIMO PRO was deliberate: the differential
kinematic model allows in-place rotation, which the navigation planner exploits to
align the vehicle heading at the start and end of each segment, and it provides a
simpler, directly verifiable odometry model (track width and wheel speed only) that is
less susceptible to unmodelled steering compliance.

The Nav2 stack consists of a global planner (NavFn/Dijkstra), which computes a
collision-free path in the pre-built occupancy grid, and a local planner (DWB), which
produces feasible velocity commands tracking that path while avoiding dynamic obstacles.
Each planner consults a layered costmap: the global costmap inflates all occupied map
cells by a fixed margin; the local costmap maintains a sensor-updated window around the
robot. A key configuration choice concerns the minimum obstacle range applied to both
costmaps. The mounted arm extends approximately 0.25--0.35\,m forward of the base in
the default rest posture; at a costmap minimum obstacle range of 0.22\,m (the default),
the arm itself is marked as a lethal obstacle in the local costmap, blocking all planned
trajectories. Raising the obstacle and raytrace minimum ranges to 0.35\,m in both
costmaps removes the arm from the observable obstacle space and restores planning. The
Gazebo differential-drive controller is configured with \texttt{odometry\_source: WORLD},
which derives odometry directly from the simulated world frame rather than integrating
wheel velocities; in simulation this yields perfect, drift-free odometry, simplifying
the validation of the navigation and manipulation pipeline before introducing realistic
odometry noise.

\section{Kinematics and Motion Planning}
The myCobot 280 is modelled as a 6-degree-of-freedom serial revolute chain. Inverse
kinematics is solved by the KDL analytical plugin, using a 22-seed random restart
strategy to increase the probability of finding a solution in the presence of
configuration-space discontinuities. Motion is planned by OMPL's RRTConnect algorithm,
which builds two trees simultaneously from the start and goal configurations and
attempts to connect them, providing fast convergence on connected planning problems.
All arm motion goals are expressed as Cartesian pose constraints on the
\texttt{gripper\_tcp} link in the \texttt{odom} planning frame; the planner computes a
collision-checked joint trajectory from the current state to the goal, which is executed
by the \texttt{ros2\_control} joint-trajectory controller.

Three named arm configurations are used by the orchestration layer: \texttt{home}, the
zero-joint-angle rest posture; \texttt{ready}, a raised forward-facing posture from
which Cartesian pick goals can be planned without joint-space discontinuities; and
\texttt{travel}, a compact downward-folded posture in which the arm's centre of mass
is positioned close to the base, reducing the inertial moment that would otherwise
excite joint oscillations during base navigation.
